% !TeX program = xelatex
\documentclass[11pt,oneside]{article}
\usepackage{ProblemsetStyle}

\hypersetup{
    pdftitle={20260902 Problem set - Solutions},
    pdfauthor={김정우},
    pdfsubject={Science Quiz mathematics practice problem set solutions},
    pdfcreator={XeLaTeX}
}

\begin{document}

\problemsettitle{20260902 Problem set}

\begin{problemsolution}{1}
수열
\[
    b_n=2\cos\left(\frac{n\pi}{4}\right)
\]
을 생각하자. 이 수열은 $a_n$과 같은 초기값을 가지며,
\[
    b_{n+2}
    =
    2\cos\left(\frac{\pi}{4}\right)b_{n+1}-b_n
    =
    \sqrt2\,b_{n+1}-b_n
\]
을 만족한다. 따라서 $a_n=b_n$이다. $2026\equiv2\pmod8$이므로
\[
    \boxed{
        a_{2026}
        =
        2\cos\left(\frac{2026\pi}{4}\right)
        =0
    }.
\]
\end{problemsolution}

\begin{problemsolution}{2}
Fibonacci sequence를 $F_0=0$, $F_1=1$로 두면 그 ordinary generating function은
\[
    \frac{1}{1-x-x^2}
    =
    \sum_{n=0}^{\infty}F_{n+1}x^n
\]
이다. 따라서 $x^{2026}$의 coefficient는 $F_{2027}$이고,
\[
    \boxed{f^{(2026)}(0)=2026!\,F_{2027}}.
\]
\end{problemsolution}

\begin{problemsolution}{3}
all-ones vector를 $\mathbf{1}$이라 하자. $\mathbf{J}_n$은 $\mathbf{1}$에 대하여 eigenvalue $n$을 가지며, $\mathbf{1}$에 수직인 모든 벡터에 대해서는 eigenvalue $0$을 갖는다. 따라서 $\mathbf{A}$의 eigenvalue는
\[
    n+1
    \quad\text{(multiplicity }1\text{)},
    \qquad
    1
    \quad\text{(multiplicity }n-1\text{)}.
\]
또한 $\mathbf{J}_n^2=n\mathbf{J}_n$이고 $\mathbf{1}$이 생성하는 직선으로의 projection matrix는 $\mathbf{J}_n/n$이므로
\[
    \mathbf{A}^{2026}
    =
    \left(\mathbf{I}_n-\frac{1}{n}\mathbf{J}_n\right)
    +(n+1)^{2026}\frac{1}{n}\mathbf{J}_n.
\]
따라서
\[
    \boxed{
        \mathbf{A}^{2026}
        =
        \mathbf{I}_n
        +
        \frac{(n+1)^{2026}-1}{n}\mathbf{J}_n
    }.
\]
\end{problemsolution}

\end{document}
